\subsection{Iteratives Design und Feedback}\label{sec:iteratives_design}

Mehrere Rückkopplungsschleifen haben die Mockups von Resteria bereits durchlaufen: Der Abgleich gegen die Konzeptbeschreibung ergänzte zuerst Challenge-Banner und Einsparangabe, später den Zugang zum Wochenplan; die Sichtung der Screens selbst korrigierte Platzierungs- und Beschriftungsfehler. Die nächste Schleife folgt vor der eigentlichen Entwicklung. In der Design-Phase (Woche 3--4, siehe~\autoref{tab:phasenplanung}) sollen Bedienführungs-Entscheidungen wie die Drei-Bereiche-Navigation aus~\autoref{sec:konzept} am Klick-Prototyp geprüft werden, bevor sie in die \ac{MVP}-Entwicklung übergehen. Auffälligkeiten lassen sich an dieser Stelle noch ohne Programmieraufwand korrigieren.

Die Feedback-Schleifen im Entwicklungsprozess von Resteria sind für diesen Projektbericht ausschließlich als geplantes Vorgehen beschrieben, nicht real durchgeführt; die Aufgabenstellung verlangt an dieser Stelle eine Beschreibung, keine tatsächliche Umsetzung. Nach einem ersten Prototyp ist eine kleine Beta-Gruppe vorgesehen, die über zwei bis drei Wochen Kernfunktionen wie das Reste-Matching und die Rettungspunkte testet. Für die Gruppengröße gibt \textcite[Kap.~17.2.2]{erlhoferWebsiteKonzeptionUndRelaunch2017} den Maßstab, da bereits wenige Testpersonen die schwerwiegenden und mehrfach auftretenden Bedienhürden einer Oberfläche offenlegen; zehn bis fünfzehn Personen sind für die Beta-Gruppe damit ausreichend dimensioniert. Rückmeldungen sollen über kurze Fragebögen und ein moderiertes Feedback-Gespräch gesammelt und in zweiwöchigen Sprints in die Weiterentwicklung einfließen. Die Fragebögen erheben je Kernaufgabe die \ac{SEQ}, eine einzelne Frage zur empfundenen Aufgabenschwierigkeit auf einer siebenstufigen Skala~\parencite[Kap.~8]{sauroLewisQuantifying2016}. Als Schwelle für eine gelungene Aufgabe legt diese Arbeit einen Mittelwert von mindestens fünf fest. Das Gespräch folgt dem Think-Aloud-Verfahren; die Testperson spricht während der Bearbeitung laut aus, was sie tut und erwartet. So werden Fehlannahmen über die Bedienführung sichtbar, die eine reine Erfolgsmessung nicht erfasst~\parencite[Kap.~6.8]{nielsenUsabilityEngineering1993}.
